\documentclass[11pt,a4paper]{article}
\usepackage[utf8]{inputenc}
\usepackage[T1]{fontenc}
\usepackage{amsfonts}
\usepackage[french]{babel}
\frenchbsetup{StandardLists=true} % à inclure si on utilise \usepackage[francais]{babel}
\usepackage{enumitem}
\usepackage{amssymb}
\def\nbR{\ensuremath{\mathrm{I\! R}}}
\usepackage{amsmath}
\usepackage{awesomebox}
\DeclareMathOperator{\e}{e}

\usepackage[left=2cm,right=2cm,top=2cm,bottom=2cm]{geometry}
\usepackage{multirow}
\usepackage[dvipsnames]{xcolor}
\usepackage{parskip}
\usepackage{caption}
\usepackage{qrcode}

\usepackage[T1]{fontenc}
\usepackage{fouriernc}
\usepackage{graphicx}
\usepackage{setspace}
\singlespacing
\usepackage{tcolorbox}
\usepackage{framed}
\usepackage{multicol}

\usepackage{fancybox}
\usepackage{fancyhdr}

\setlength{\headheight}{14pt}
\pagestyle{fancy}
\renewcommand\headrulewidth{1pt}
\fancyhead[L]{Lycée Denis Diderot }
\fancyhead[R]{2BTS}
\setcounter{page}{1}\fancyfoot[C]{page \thepage}
\fancyfoot[R]{\today}
\fancyfoot[L]{Alexis RUIZ}
\usepackage{tikz}
\usepackage{pgfplots}
\pgfplotsset{compat=1.18}
\usepackage{tkz-tab}
\usetikzlibrary{arrows}
\usepackage{pifont}

\usepackage{tipfr}

\newcolumntype{C}[1]{>{\centering\arraybackslash}p{#1cm}}
\renewcommand{\arraystretch}{1.3} 

\newcommand{\ud}{\mathrm{d}}
\newcommand{\V}{\overrightarrow}
\newcommand{\Rep}{(O;\V{u};\V{v})} 
\newcommand{\dd}{
		\mathop{}\mathopen{}\mathrm{d}
	}
\newcommand{\defbox}{\awesomebox[black]{2pt}{\faEye}{black}}
\newcounter{numexos}%Création d'un compteur qui s'appelle numexos
\setcounter{numexos}{0}%initialisation du compteur
\newcommand{\exercice}[1]{%Création d'une macro ayant un paramètre
\addtocounter{numexos}{1}%chaque fois que cette macro est appelée, elle ajoute 1 au compteur numexos
\textcolor{black}{Exercice\,\thenumexos\,}\,#1%Met en noir Exercice et la valeur du compteur appelée par \thenumeexos
}

\newcounter{numact}%Création d'un compteur qui s'appelle numexos
\setcounter{numact}{0}%initialisation du compteur
\newcommand{\activit}[1]{%Création d'une macro ayant un paramètre
\addtocounter{numact}{1}%chaque fois que cette macro est appelée, elle ajoute 1 au compteur numexos
\textcolor{black}{A\thenumexos\,}\,#1%Met en noir Exercice et la valeur du compteur appelée par \thenumeexos
}

\newcounter{nump}%Création d'un compteur qui s'appelle numexos
\setcounter{nump}{0}%initialisation du compteur
\newcommand{\propriete}[1]{%Création d'une macro ayant un paramètre
\addtocounter{nump}{1}%chaque fois que cette macro est appelée, elle ajoute 1 au compteur numexos
\textcolor{black}{Propriété\,\thenump\,}\,#1%Met en noir Exercice et la valeur du compteur appelée par \thenumeexos
}



\newcommand{\bex}[2]{%
\def\bextitle{#1}%
\addtocounter{numexos}{1}%
\begin{tcolorbox}[colback=white,colframe=black!30,arc=3pt,boxrule=0.6pt,
  left=4mm,right=4mm,top=1mm,bottom=2mm]%
  \textbf{Exercice~\thenumexos\ifx\bextitle\empty\else~:~~#1\fi}\par
  #2%
\end{tcolorbox}}

\newcommand{\act}[3]
{\begin{bclogo}[couleur = white,logo=,  arrondi = 0.1,barre=none]
{~\activit:~~{#2}}
~~
{#2}
~~
~~
\end{bclogo}
\vspace{3mm}
\cadre{{#3}}
}

\newcommand{\rem}[1]
{\begin{bclogo}[couleur = white,logo=\bcinfo , arrondi = 0.1]
{~Remarque~:}
~~
{#1}
~~
\end{bclogo}}


\newcommand{\bthe}[1]
{\begin{bclogo}[couleur = white,logo=\bcinfo , arrondi = 0.1]
{~Théorème~:}
~~
{#1}
~~
\end{bclogo}}



\newcommand{\bpre}[1]
{\begin{bclogo}[couleur = white,logo=, barre=none, arrondi = 0.1]
{~Prérequis~:}
~~
{#1}
~~
\end{bclogo}}



\newcommand{\intro}[1]
{\begin{bclogo}[couleur = white,logo=\bcinfo , arrondi = 0.1]
{~Introduction~:}
~~
{#1}
~~
\end{bclogo}}


\newcommand{\bdef}[1]
{\begin{bclogo}[couleur = white,logo=\bcinfo , arrondi = 0.1]
{~Définition~:}
~~
{#1}
~~
\end{bclogo}}


\newcommand{\bprop}[2]
{\begin{bclogo}[couleur =white ,logo= , barre=none,  arrondi = 0.1]
{~\propriete:~~{#1}}
~~
{#2}
~~
\end{bclogo}}




\newcommand{\cadre}[1]
{\begin{bclogo}[couleur = white,logo=,  arrondi = 0.1,barre=none]
~~
{%
\rule{0mm}{#1mm}}\par
\medskip
~~
\end{bclogo}}


\newcommand{\cadreblanc}[1]{%
\medskip\framebox[\linewidth][c]{%
\rule{0mm}{#1mm}}\par
\medskip}

\setlength{\columnseprule}{.25mm}
\usepackage{multirow,tabularx,booktabs}
\usepackage{colortbl}
\usepackage{wrapfig}
\usepackage[np]{numprint}

\usepackage[tikz]{bclogo} 

\def\N{{\mathbb N}}
\def\Z{{\mathbb Z}}
\def\R{{\mathbb R}}
\def\C{{\mathbb C}}
\def\e{{\text{e}}}
\def\dt{{\text{d}t}}

%%%%%%%%%%%%%%%%%%%%%%%%%%%%%%%%%%%%%%%%%%%%%%%%%%%%%%%%%%%%%%%%%%%%%%%%%%%%%%
%%%%%%%%%%%%%%%%%%%%%%%%%%%%%%%%%%%%%%%%%%%%%%%%%%%%%%%%%%%%%%%%%%%%%%%%%%%%%%
%%%%%%%%%%%%%%%%%%%%%%%%%%%%%%%%%%%%%%%%%%%%%%%%%%%%%%%%%%%%%%%%%%%%%%%%%%%%%%


\begin{document}





\newtcolorbox{mybox}{colback=white!10!white,
colframe=black!5!black}
\begin{mybox}
\begin{center}    
\textbf{\LARGE CHAPITRE 10 - Séries de Fourier}\end{center}
\end{mybox}

\vfill


\vspace{0.1cm}
\begin{center}
\begin{minipage}{14cm}

    \begin{bclogo}[couleur = white,  arrondi = 0.1]
{~Capacités attendues~:}
~~
\begin{itemize}
\item Calculer les coefficients de Fourier dans des cas simples.
\item Savoir identifier le bon développement en sérier de Fourier d'une fonction donnée. 
\item Savoir calculer la valeur approchée d'une valeur efficace d'une fonction en utilisant son développement en série de Fourier.

    
\end{itemize}

\end{bclogo}
\end{minipage}
    
\end{center}
  

\bpre{%
\setstretch{0.85}
On considère la fonction $f$ définie sur $\mathbb{R}$ par les conditions suivantes : \\
\begin{center}
$\begin{cases}
f(x) &\textrm{est paire}\\
f(x)&\textrm {est de période }2\pi \\
f(x)=1  & \textrm {pour~~} x \in \left[0;\dfrac{\pi}{2} \right]\\
f(x)=0  & \textrm {pour~~} x \in \left]\dfrac{\pi}{2};\pi \right]\\
\end{cases}$
\end{center}

\begin{enumerate}
    \item Représenter dans le repère ci-dessous la fonction $f$ sur l'intervalle $[-\pi;2\pi]$.
\end{enumerate}
}

\begin{center}

\begin{tikzpicture}[x=2cm, y=4cm]
\def\xmin{-2.3} \def\xmax{4.3} \def\ymin{-0.15} \def\ymax{1.25}
\draw[xstep=0.25,ystep=0.1,very thin,orange!25] (\xmin,\ymin) grid (\xmax,\ymax);
\draw[xstep=1,ystep=0.5,thin,orange!55] (\xmin,\ymin) grid (\xmax,\ymax);
\draw[->,>=stealth',thick,blue!70!black] (\xmin,0)--(\xmax,0) node[right]{$x$};
\draw[->,>=stealth',thick,blue!70!black] (0,\ymin)--(0,\ymax) node[above]{$y$};
\foreach \n/\lab in {-2/{-\pi},-1/{-\frac{\pi}{2}},1/{\frac{\pi}{2}},2/{\pi},3/{\frac{3\pi}{2}},4/{2\pi}}
    \draw[blue!70!black] (\n,2pt)--(\n,-6pt) node[below]{\small$\lab$};
\foreach \y in {0.5,1}
    \draw[blue!70!black] (2pt,\y)--(-5pt,\y) node[left]{\small$\y$};
\draw[black!60,thick] (\xmin,\ymin) rectangle (\xmax,\ymax);
\end{tikzpicture}

\end{center}

\bpre{%
\setstretch{0.85}
\begin{enumerate}
\setcounter{enumi}{1}
\item Tracer la courbe représentative de la fonction $g(x)=\dfrac{1}{2}+\dfrac{2}{\pi}\cos(x)-\dfrac{2}{3\pi}\cos(3x)$ \\[-2mm]
\item Tracer la courbe représentative de la fonction $h(x)=g(x)+\dfrac{2}{5\pi}\cos(5x)$ \\[-2mm]
\item  Tracer la courbe représentative de la fonction $k(x)=h(x)-\dfrac{2}{7\pi}\cos(7x)$\\[-2mm]
\item Conjecturer une fonction dont la représentation graphique s'approcherait encore plus de celle de la fonction $f$.\\[-2mm]

\dotfill

\end{enumerate}




}

\vfill 

%%%%%%%%%%%%%%%%%%%%%%%%%%%
\newpage
%%%%%%%%%%%%%%%%%%%%%%%%%%%

\begin{mybox}
    \textbf{1. Calculs des coefficients de Fourier}
\end{mybox}

\begin{mybox}
    \textbf{1.1. Les séries de Fourier}
\end{mybox}

\bdef{
\vspace{5cm}
}

Les nombres $a_0$, $a_n$ et $b_n$ \textbf{sont les coefficients de Fourier } de la fonction $f$.

\textbf{Exemple :} Dans le pré requis, vous retrouvez le début du développement d'une série de Fourier d'une fonction $f$. \\

Retrouver les coefficients suivants : $a_0=$ ~~~~~~ , $a_1=$ ~~~~~~ , $a_2=$ ~~~~~~, $a_3=$ ~~~~~~, $a_4=$ ~~~~~~, $a_5=$ ~~ \\[3mm]

Quelle est la valeur des coefficient $b_n$ : \dotfill

\rem{ Si $T$ (en seconde , $s$) est la période d'une fonction $f$ alors $\omega$ est appelé sa pulsation (en radians par seconde , $rad.s^{-1}$) , on note alors la fréquence d'un signal $f=\dfrac{\omega}{2\pi}$. Son unité étant le Hertz ($Hz$).

}

\begin{mybox}
    \textbf{2.2. Calculs des coefficients}
\end{mybox}


Nous admettrons les résultats suivants : \fbox{$a_0=\dfrac{1}{T} \displaystyle \int_\alpha^{\alpha+T} f(t)\mathrm{d}t$}\\[3mm]
$a_0$ est en fait la \textbf{valeur moyenne} de la fonction $f$ sur une période.
\vspace{3mm}

\rem{Le réel $\alpha$ est choisi arbitrairement et le résultat de l'intégrale ne dépend pas de ce choix. On prendra le plus souvent $\alpha = 0$ ou $\alpha = \dfrac{-T}{2}.$

}


\vspace{0.3cm}

\begin{multicols}{2}
\begin{center}
    

	\fbox{$a_n=\dfrac{2}{T} \displaystyle \int_\alpha^{\alpha+T} f(t)\cos(n\omega t)\mathrm{d}t$} \\
	\fbox{$b_n=\dfrac{2}{T} \displaystyle \int_\alpha^{\alpha+T} f(t)\sin(n\omega t)\mathrm{d}t$} \\
\end{center}\end{multicols}

\vspace{0.5cm}

\rem{Si la fonction $f$ est de période $2\pi$ alors on a $\omega = 1$ et :\\[3 mm]
$a_0=\dfrac{1}{2\pi} \displaystyle \int_\alpha^{\alpha+2\pi} f(t)\mathrm{d}t$,\hfill $a_n=\dfrac{1}{\pi} \displaystyle \int_\alpha^{\alpha+2\pi} f(t)\cos(n t)\mathrm{d}t$, \hfill $b_n=\dfrac{1}{\pi} \displaystyle \int_\alpha^{\alpha+2\pi} f(t)\sin(n t)\mathrm{d}t$.

}


\begin{bclogo}[couleur = white,logo=,  arrondi = 0.1,barre=none]
{Activité 1. Calculer les coefficients de Fourier à la main.}
~
\end{bclogo}


On considère la fonction définie par : 
$\begin{cases}
   	f(x)&\textrm {est de période }2\pi \\
	f(x)=1  & \textrm {pour~~} x \in \left[0;\pi \right]\\
	f(x)=0  & \textrm {pour~~} x \in \left]\pi;2\pi \right]\\
\end{cases}$

\vfill 

\begin{center}
\fbox{
\begin{minipage}{16cm}{
\textbf{Calcul de $\mathbf{a_0}$ : }\\
\vspace{1.5cm}
}

\end{minipage}
}
\vfill 
\fbox{
\begin{minipage}{16cm}{
\textbf{Calcul de $\mathbf{a_n}$ : }\\
\vspace{2cm}
}

\end{minipage}
}
\vfill 
\fbox{
\begin{minipage}{16cm}{
\textbf{Calcul de $\mathbf{b_n}$ : }\\
\vspace{3cm}
}

\end{minipage}
}
\vfill 
\fbox{
\begin{minipage}{16cm}{
\textbf{Le développement en série de Fourier de la fonction $\mathbf{f}$ est donc : }\\
\vspace{1cm}
}

\end{minipage}
}
\vfill 
\end{center}


%%%%%%%%%%%%%%%%%%%%%%%%%%%%%%%%%%%%%%%%%%%%%%%%%%%%
\newpage
%%%%%%%%%%%%%%%%%%%%%%%%%%%%%%%%%%%%%%%%%%%%%%%%%%%%


%%%%%%%%%%%%%%%%%%%%%%%%%%%%%%%%%%%%%%%%%%%%%%%%%%%%%%%%%%
\newpage
%%%%%%%%%%%%%%%%%%%%%%%%%%%%%%%%%%%%%%%%%%%%%%%%%%%%%%%%%%




\bex{}{}

On considère la fonction définie par : 
$\begin{cases}
   	f(x)&\textrm {est de période }2\pi \\
	f(x)=1  & \textrm {pour~~} x \in \left[0;\dfrac{\pi}{2} \right]\\
	f(x)=0  & \textrm {pour~~} x \in \left]\dfrac{\pi}{2};2\pi \right]\\
\end{cases}$

calculer ses coefficients de Fourier. 
\begin{center}
\fbox{
\begin{minipage}{16cm}{
\textbf{Calcul de $\mathbf{a_0}$ : }\\
\vspace{1cm}
}

\end{minipage}
}

\fbox{
\begin{minipage}{16cm}{
\textbf{Calcul de $\mathbf{a_n}$ : }\\
\vspace{3cm}
}

\end{minipage}
}

Calculer les coefficients de $n=1$ à $n=5$
\begin {multicols}{5}
	\cadreblanc{8}
	\cadreblanc{8}
	\cadreblanc{8}
	\cadreblanc{8}
	\cadreblanc{8}
\end{multicols}

\end{center}

\fbox{
\begin{minipage}{16cm}{
\textbf{Calcul de $\mathbf{b_n}$ : }\\
\vspace{3cm}
}

\end{minipage}
}

Calculer les coefficients de $n=1$ à $n=5$
\begin {multicols}{5}
	\cadreblanc{8}
	\cadreblanc{8}
	\cadreblanc{8}
	\cadreblanc{8}
	\cadreblanc{8}
\end{multicols}

\fbox{
\begin{minipage}{16cm}{
\textbf{Écrire le développement en série de Fourier de la fonction $f$ pour jusqu'à $\mathbf{n=4}$}\\
\vspace{3cm}
}

\end{minipage}
}




\newpage

\bex{}{}

On considère la fonction définie par : 
$\begin{cases}
   	f(x)&\textrm {est de période }2\pi \\
	f(x)=1  & \textrm {pour~~} x \in \left[0;\pi \right]\\
	f(x)=-1  & \textrm {pour~~} x \in \left]\pi;2\pi \right]\\
\end{cases}$

calculer ses coefficients de Fourier. 
\begin{center}
\fbox{
\begin{minipage}{16cm}{
\textbf{Calcul de $\mathbf{a_0}$ : }\\
\vspace{1cm}
}

\end{minipage}
}

\fbox{
\begin{minipage}{16cm}{
\textbf{Calcul de $\mathbf{a_n}$ : }\\
\vspace{3cm}
}

\end{minipage}
}

Calculer les coefficients de $n=1$ à $n=5$
\begin {multicols}{5}
	\cadreblanc{8}
	\cadreblanc{8}
	\cadreblanc{8}
	\cadreblanc{8}
	\cadreblanc{8}
\end{multicols}


\fbox{
\begin{minipage}{16cm}{
\textbf{Calcul de $\mathbf{b_n}$ : }\\
\vspace{3cm}
}

\end{minipage}
}

Calculer les coefficients de $n=1$ à $n=5$
\begin {multicols}{5}
	\cadreblanc{8}
	\cadreblanc{8}
	\cadreblanc{8}
	\cadreblanc{8}
	\cadreblanc{8}
\end{multicols}

\fbox{
\begin{minipage}{16cm}{
\textbf{Écrire le développement en série de Fourier de la fonction $f$ pour jusqu'à $\mathbf{n=4}$}\\
\vspace{3cm}
}

\end{minipage}
}
\end{center}
\newpage

\bex{}{}

On considère la fonction définie par : 
$\begin{cases}
   	f(x)&\textrm {est de période }2 \\
	f(x)=1  & \textrm {pour~~} 0\le x <1 \\
	f(x)=0  & \textrm {pour~~} 1\le x <2 \\
\end{cases}$

calculer ses coefficients de Fourier. 
\begin{center}
\fbox{
\begin{minipage}{16cm}{
\textbf{Calcul de $\mathbf{a_0}$ : }\\
\vspace{1cm}
}

\end{minipage}
}

\fbox{
\begin{minipage}{16cm}{
\textbf{Calcul de $\mathbf{a_n}$ : }\\
\vspace{3cm}
}

\end{minipage}
}

Calculer les coefficients de $n=1$ à $n=5$
\begin {multicols}{5}
	\cadreblanc{8}
	\cadreblanc{8}
	\cadreblanc{8}
	\cadreblanc{8}
	\cadreblanc{8}
\end{multicols}



\fbox{
\begin{minipage}{16cm}{
\textbf{Calcul de $\mathbf{b_n}$ : }\\
\vspace{3cm}
}

\end{minipage}
}

Calculer les coefficients de $n=1$ à $n=5$
\begin {multicols}{5}
	\cadreblanc{8}
	\cadreblanc{8}
	\cadreblanc{8}
	\cadreblanc{8}
	\cadreblanc{8}
\end{multicols}
\end{center}
\fbox{
\begin{minipage}{16cm}{
\textbf{Écrire le développement en série de Fourier de la fonction $f$ pour jusqu'à $\mathbf{n=4}$}\\
\vspace{3cm}
}

\end{minipage}
}

%%%%%%%%%%%%%%%%%%%%%%%%%%%
\newpage
%%%%%%%%%%%%%%%%%%%%%%%%%%

\begin{mybox}
    3. Calcul des coefficients de Fourier dans les cas particuliers 
\end{mybox}



\begin {center}
\begin{minipage}{14cm}{

\begin{bclogo}[couleur = white,logo=\bcetoile,  arrondi = 0.1]
{Bon à savoir} 

\vspace{3cm}

\end{bclogo}

}

\end{minipage}

\end{center}

\vfill

\begin{mybox}
    3.1. Si la fonction $f$ est impaire
\end{mybox}

\vspace{3mm}
\cadre{30}

\vfill

\begin{mybox}
   3.2. Si la fonction $f$ est paire
\end{mybox}
\vspace{3mm}
\cadre{40}

\vfill 



\begin {center}
\begin{minipage}{14cm}{

\begin{bclogo}[couleur = white,logo=\bcetoile,  arrondi = 0.1]
{Bon à savoir} 

\vspace{3cm}

\end{bclogo}

}

\end{minipage}

\end{center}

\vfill 


\newpage

\fcolorbox[gray]{0}{0.9}{\textbf{A4. Calculer les coefficients de Fourier pour une fonction paire.}}\\[5mm]
On considère la fonction $f$ de période $\pi$, paire et définie par $f(x)=x$ sur l'intervalle $\left[0;\dfrac{\pi}{2}\right]$\\

Dans ce genre d'exercice il est préférable de commencer par représenter la fonction sur l'intervalle permettant de voir la périodicité et la parité de la fonction.

\begin{center}
    
\begin{tikzpicture}[x=1.4cm, y=2.2cm]
\def\xmin{-4.4} \def\xmax{4.4} \def\ymin{-1.4} \def\ymax{1.4}
\draw[xstep=0.25,ystep=0.1,very thin,orange!25] (\xmin,\ymin) grid (\xmax,\ymax);
\draw[xstep=1,ystep=0.5,thin,orange!55] (\xmin,\ymin) grid (\xmax,\ymax);
\draw[->,>=stealth',thick,blue!70!black] (\xmin,0)--(\xmax,0) node[right]{$x$};
\draw[->,>=stealth',thick,blue!70!black] (0,\ymin)--(0,\ymax) node[above]{$y$};
\foreach \n/\lab in {-4/{-2\pi},-3/{-\frac{3\pi}{2}},-2/{-\pi},-1/{-\frac{\pi}{2}},1/{\frac{\pi}{2}},2/{\pi},3/{\frac{3\pi}{2}},4/{2\pi}}
	\draw[blue!70!black] (\n,2pt)--(\n,-6pt) node[below]{\small$\lab$};
\foreach \y in {-1,-0.5,0.5,1}
	\draw[blue!70!black] (2pt,\y)--(-5pt,\y) node[left]{\small$\y$};
\draw[black!60,thick] (\xmin,\ymin) rectangle (\xmax,\ymax);
\end{tikzpicture}
\end{center}

\cadreblanc{120}

%%%%%%%%%%%%%%%%%%%%%%%%%%%%%%%%%%%%%%%%
\newpage
%%%%%%%%%%%%%%%%%%%%%%%%%%%%%%%%%%%%%%%%


\bex{Faire le point avec un QCM}{
Pour chaque question indiquer la bonne réponse. \\
Pour chacune des question, on note $$S(t)=a_0+ \sum\limits_{n=1}^{+\infty} (a_n\cos(n\omega x)+b_n\sin(n\omega x))$$ le développement en série de Fourier de la fonction $f$.  

\bigskip

1. Si pour tout entier naturel $n > 0$, $a_n=0$ alors la fonction $f$ :
\begin{itemize}[label=$\square$]
  \item   a) est paire
  \item   b) est impaire
  \item   c) on ne peut pas savoir
\end{itemize}
\bigskip
2. Si la fonction $f$ est paire, périodique de période $2\pi$ et si sur $\left[0;\dfrac{\pi}{2}\right[ f(x) = \pi$ et sur $\left[\dfrac{\pi}{2}; \pi\right[ f(x) = 0$ alors pour tout entier naturel non nul $n$ on a : 
\begin{itemize}[label=$\square$]
\item   a) $a_n=\dfrac{1}{\pi}\displaystyle \int_0^{\pi} \pi \cos(nt) \mathrm{d}t$
\item   b) $a_n=\dfrac{2}{\pi}\displaystyle \int_0^{\frac{\pi}{2}} \pi \cos(nt) \mathrm{d}t$
\item   c) $a_n=\dfrac{2}{\pi}\displaystyle \int_0^{\pi} \pi \cos(nt) \mathrm{d}t$
\end{itemize}
\bigskip
3. Si la fonction $f$ est  périodique de période $\pi$ et si sur $\left[0;\dfrac{\pi}{2}\right[ f(x) = \pi$ et sur $\left[\dfrac{\pi}{2}; \pi\right[ f(x) = 0$ alors pour tout entier naturel non nul $n$ on a : 
\begin{itemize}[label=$\square$]
\item   a) $a_0=\pi$ et pour $n \ne 0$ on a $a_n=0$
\item   b) pour $n \ne 0$ on a $b_n=0$ 
\item   c) Aucun des coefficients n'est nul. 
\end{itemize}
\bigskip
4. Si la fonction $f$ est  périodique de période $2\pi$ et si sur $[0;2\pi[ f(x) = x^2$ alors : 
\begin{itemize}[label=$\square$]
\item   a) $f$ est paire
\item   b) $f$ est impaire
\item   c) $f$ est ni paire ni impaire
\end{itemize}
\bigskip
5. Si une fonction a pour période $T=4$ alors $\omega = $
\begin{itemize}[label=$\square$]
  \item   a) $\dfrac{\pi}{2}$
  \item   b) on ne peut pas savoir
  \item   c) $\dfrac{1}{4}$
  
\end{itemize}
}
\bigskip

\newpage

\bex{}{}
On considère la fonction $f$ définie sur $\mathbb{R}$, impaire et périodique de période $\pi$ telle que :  
$$\begin{cases}
	f(x)=\dfrac{1}{2}  & \textrm {si~~} x \in \left[0;\dfrac{\pi}{4} \right[\\[5mm]
	f(x)=0  & \textrm {si~~} x \in \left]\dfrac{\pi}{4};\dfrac{\pi}{2} \right]\\
\end{cases}$$

Calculer ses coefficients de Fourier. 
\begin{center}
\fbox{
\begin{minipage}{16cm}{
\textbf{Calcul de $\mathbf{a_0}$ : }\\
\vspace{1cm}
}
\vspace{1 cm}
\end{minipage}
}
\end{center}
\vspace{1 cm}
\begin{center}
\fbox{
\begin{minipage}{16cm}{
\textbf{Calcul de $\mathbf{a_n}$ : }\\
\vspace{2cm}
}

\end{minipage}
}
\end{center}
\vspace{1 cm}
\begin{center}
\fbox{
\begin{minipage}{16cm}{
\textbf{Calcul de $\mathbf{b_n}$ : }\\
\vspace{7cm}
}

\end{minipage}
}
\end{center}
\vspace{1 cm}
\begin{center}
\fbox{
\begin{minipage}{16cm}{
\textbf{Calcul de $S(x)$ : }\\
\vspace{1cm}
}

\end{minipage}
}
\end{center}


\newpage
%%%%%%%%%%%%%%%%%%%%%%%%%%%%%%


\fcolorbox[gray]{0}{0.9}{\textbf{4. La formule de Parseval et ses applications}}\\[3 mm]

\fcolorbox[gray]{0}{0.9}{\textbf{4.1. La valeur efficace d'un signal.}}\\[3 mm]

\begin{center}
	\fbox{
			Si $f$ est une fonction périodique de période $T$, sa valeur efficace est : $f_{eff}(x)=\sqrt{\dfrac{1}{T}\displaystyle \int_\alpha^{\alpha+T} (f(x))^2\mathrm{d}x}$
			}

\end{center}
\vspace{5mm}

Cette formule est souvent utilisée en sciences physiques pour calculer la tension efficace, l'intensité efficace ... \\[5 mm]

\fcolorbox[gray]{0}{0.9}{\textbf{4.1. La valeur efficace d'une fonction développable en série de Fourrier.}}\\[3 mm]
Considérons une fonction de période $T$ dont le développement en série de Fourier est : $$f(x)=a_0+ \sum\limits_{n=1}^{+\infty}\left(a_n\cos(n\omega x) + b_n \sin (n\omega x)\right)$$

Nous admettrons \textbf{la formule de Parceval} : \\

\begin{center}
\fbox{$\dfrac{1}{T}\displaystyle \int_\alpha^{\alpha+T} (f(x))^2\mathrm{d}x=a_0^2+\dfrac{1}{2} \displaystyle\sum\limits_{n=1}^{+\infty}\left(a_n^2+b_n^2\right)$}

\end{center}
\vspace{5 mm}
Si on note $f_{eff}$ la valeur efficace de la fonction $f$ ce résultat peut aussi s'écrire :

\begin{center}
\fbox{$f_{eff}^2=a_0^2+\dfrac{1}{2} \displaystyle\sum\limits_{n=1}^{+\infty}\left(a_n^2+b_n^2\right)$}
\end{center}

\newpage

\textbf{A5. Valeur efficace d'une fonction}

On considère la fonction définie par : 
$\begin{cases}
    f(x)&\textrm  {est paire}\\
   	f(x)&\textrm {est de période }2\pi \\
	f(x)=1  & \textrm {pour~~} x \in \left[0;\dfrac{\pi}{2} \right]\\
	f(x)=0  & \textrm {pour~~} x \in \left]\dfrac{\pi}{2};\pi \right]\\
\end{cases}$

Représenter $f$ dans le repère ci-dessous
\begin{center}
\begin{tikzpicture}[x=1.2cm, y=2.2cm][scale=0.8]
\def\xmin{-4.4} \def\xmax{6.4} \def\ymin{-1.4} \def\ymax{1.4}
\draw[xstep=0.25,ystep=0.1,very thin,orange!25] (\xmin,\ymin) grid (\xmax,\ymax);
\draw[xstep=1,ystep=0.5,thin,orange!55] (\xmin,\ymin) grid (\xmax,\ymax);
\draw[->,>=stealth',thick,blue!70!black] (\xmin,0)--(\xmax,0) node[right]{$x$};
\draw[->,>=stealth',thick,blue!70!black] (0,\ymin)--(0,\ymax) node[above]{$y$};
\foreach \n/\lab in {-4/{-2\pi},-3/{-\frac{3\pi}{2}},-2/{-\pi},-1/{-\frac{\pi}{2}},1/{\frac{\pi}{2}},2/{\pi},3/{\frac{3\pi}{2}},4/{2\pi},5/{\frac{5\pi}{2}},6/{3\pi}}
	\draw[blue!70!black] (\n,2pt)--(\n,-6pt) node[below]{\small$\lab$};
\foreach \y in {-1,-0.5,0.5,1}
	\draw[blue!70!black] (2pt,\y)--(-5pt,\y) node[left]{\small$\y$};
\draw[black!60,thick] (\xmin,\ymin) rectangle (\xmax,\ymax);
\end{tikzpicture}
\end{center}

Nous avons obtenu le début du développement en série de Fourier de la fonction $f$ donné par $$g(x)=\dfrac{1}{2}+\dfrac{2}{\pi}\cos(x)-\dfrac{2}{3\pi}\cos(3x)$$
Calculer les valeurs efficaces $f_{eff}$ et $g_{eff}$ des deux fonctions $f$ et $g$ et donner la valeur approchée de leur quotient.\\ 
\cadreblanc{80}

\textbf{4.3 Applications}\\
En physique cette formule permet de voir si l'utilisation des premiers harmoniques permet d'utiliser le début du développement en série de Fourier pour remplacer un signal \textbf{(voir A4)}. 

%%%%%%%%%%%%%%%%%%%%%%%%%%%%%%%
\newpage
%%%%%%%%%%%%%%%%%%%%%%%%%%%%%%%%

\large
\begin{center} 
\textbf {FICHE ESSENTIEL}\\[0.5cm]
\normalsize

\end{center}
\hrule
\vspace{8 mm}
\begin{center}
\fbox{
\begin{minipage}{17cm}{
\textbf{Les séries numériques}\\

Les séries géométriques :\\
La série $\displaystyle \sum\limits_{n=0}^{+\infty}q^n$ converge si et seulement si $q \in ]-1;1[$.\\

Les séries de Riemann :\\
La série $\displaystyle \sum\limits_{n=1}^{+\infty}\dfrac{1}{n^\alpha}$ converge si et seulement si $n>1$.\\[2mm]

\vspace{8 mm}
\textbf{Le développement en série de Fourier d'une fonction}\\
Si $f$ est une fonction périodique de période $T$, on note \textbf{sa pulsation} $\omega = \dfrac{2\pi}{T}$ \\ 
Son développement en série de Fourier est $S(x)=a_0+\displaystyle \sum\limits_{n=1}^{+\infty}(a_n\cos(n\omega x) + b_n\sin(n\omega x)$ avec :\\
\begin{center}
$a_0=\dfrac{1}{T}\displaystyle \int\limits_{\alpha}^{\alpha+T}f(x)~~\mathrm{d}x$~~~~ il s'agit de la valeur moyenne de $f$.\\
$a_n=\dfrac{2}{T}\displaystyle \int\limits_{\alpha}^{\alpha+T}f(x)\cos(n\omega x)~~\mathrm{d}x$~~~~ et~~~~ $b_n=\dfrac{2}{T}\displaystyle \int\limits_{\alpha}^{\alpha+T}f(x)\sin(n\omega x)~~\mathrm{d}x$
\end{center}

\vspace{8 mm}
\textbf{Cas particuliers}\\

Si $f$ est une fonction impaire\\ 
tous les coefficients $a_n$ y compris $a_0$ sont nuls et\\  \begin{center}
\fbox{$b_n=\dfrac{4}{T}\displaystyle \int\limits_{0}^{\frac{T}{2}}f(x)\sin(n\omega x)~~\mathrm{d}x$}
\end{center}

Si $f$ est une fonction paire\\ 
tous les coefficients $b_n$ sont nuls et\\ \begin{center}
\fbox{$a_n=\dfrac{4}{T}\displaystyle \int\limits_{0}^{\frac{T}{2}}f(x)\cos(n\omega x)~~\mathrm{d}x$}
\end{center} 

\vspace{8 mm}

\textbf{La formule de Parseval}\\
Si la fonction $f$ de période $T$ est développable ne série de Fourier sa valeur efficace est donnée en fonction des coefficients de Fourier par la formule : \\
\begin{center}
\fbox{$f_{eff}^2=\dfrac{1}{T}\displaystyle \int\limits_{\alpha}^{\alpha+T}(f(x))^2)~~\mathrm{d}x=a_0^2+\dfrac{1}{2}\displaystyle \sum\limits_{n=1}^{+\infty}(a_n^2 + b_n^2)$} 
\end{center}




\vspace{3 mm}
}

\end{minipage}}
\end{center}

%%%%%%%%%%%%%%%%%%%%%%%%%%%%%%%%%%%%%%
\newpage
%%%%%%%%%%%%%%%%%%%%%%%%%%%%%%%%%%%%%%
\fcolorbox[gray]{0}{0.9}{\textbf{Extrait - Brevet de technicien supérieur 10 mai 2021 - groupement A}}\\[2mm]
\textbf{Exercice 1 \hfill 11 points}

\medskip

La communication par courants porteurs en ligne (ou CPL) permet de transmettre des informations en utilisant des conducteurs électriques en fonctionnement.

Le principe des CPL consiste à superposer au réseau électrique un signal de haute fréquence et de basse énergie. Ce deuxième signal se propage sur l'installation électrique et peut être reçu et décodé à distance.

\textbf{Partie A : Étude d'un signal}

\medskip

On s'intéresse à une tension périodique $U$. On note $T$ sa période.

\medskip

\begin{enumerate}
\item On donne ci-dessous une représentation graphique de $U$, exprimée en volt (V), en fonction de $t$ \textbf{exprimé en microseconde} ($\mu s$). Rappel: $1\mu s = 10^{-6}~s$.

\begin{center}
\begin{tikzpicture}[x=0.45cm, y=0.35cm]
% Axes
\def\xmin{-11} \def\xmax{20} \def\ymin{-1.5} \def\ymax{14}
\draw[xstep=1,ystep=1,very thin,orange!25] (\xmin,\ymin) grid (\xmax,\ymax);
\draw[xstep=2,ystep=4,thin,orange!55] (\xmin,\ymin) grid (\xmax,\ymax);
\draw[->,>=stealth',thick,blue!70!black] (\xmin,0)--(\xmax,0) node[right]{\small$t~(\mu s)$};
\draw[->,>=stealth',thick,blue!70!black] (0,\ymin)--(0,\ymax) node[above]{\small$U~(\text{V})$};
% Graduation x
\foreach \x in {-10,-8,-6,-4,-2,2,4,6,8,10,12,14,16,18}
	\draw[blue!70!black] (\x,2pt)--(\x,-5pt) node[below]{\tiny$\x$};
% Graduation y
\foreach \y in {4,8,12}
	\draw[blue!70!black] (2pt,\y)--(-5pt,\y) node[left]{\small$\y$};
% Signal : période 10, haut 12V pendant 8μs, bas 0 pendant 2μs
% Périodes visibles : [-10,0), [0,10), [10,20)
% Période [-10, 0) : U=12 sur [-10,-2], U=0 sur [-2,0]
\draw[very thick,red!70!black] (-10,12)--(-2,12)--(-2,0)--(0,0);
% Période [0, 10) : U=12 sur [0,8], U=0 sur [8,10]
\draw[very thick,red!70!black] (0,12)--(8,12)--(8,0)--(10,0);
% Période [10, 20) : U=12 sur [10,18], U=0 sur [18,19]
\draw[very thick,red!70!black] (10,12)--(18,12)--(18,0)--(19,0);
% Traits verticaux (discontinuités)
\draw[very thick,red!70!black] (-10,0)--(-10,12);
\draw[very thick,red!70!black] (0,0)--(0,12);
\draw[very thick,red!70!black] (10,0)--(10,12);
\draw[black!60,thick] (\xmin,\ymin) rectangle (\xmax,\ymax);
\end{tikzpicture}
\end{center}

	\begin{enumerate}
		\item D'après la représentation graphique ci-dessus, quelle est la valeur de $T$ en microseconde ?\\
		\cadreblanc{10}
		\item La fréquence d'un signal périodique, en Hz, suit la formule $f =  \dfrac{1}{T}$ où $T$ est exprimé en seconde. 
		
Quelle est la fréquence de la tension périodique $U$ ?\\
\cadreblanc{10}
	\end{enumerate}
\item On modélise l'évolution de $U$ (en volt) en fonction de $t$ (en microseconde) à l'aide d'une fonction numérique $f$ définie sur $\mathbb{R}$. Ainsi : $U = f(t)$.

On admet que la fonction $f$ est paire, périodique de période $T$, développable en série de Fourier et vérifie, pour tout réel $t$ :

\[f(t) = a_0 + \displaystyle\sum_{n=1}^{+ \infty} \left(a_n \cos (n \omega t) + b_n \sin(n\omega t) \right)\quad \text{avec}\: \omega = \dfrac{2\pi}{T}.\]

	\begin{enumerate}
		\item Quelle est la valeur de $b_n$ pour tout entier $n \geqslant 1$ ? Justifier. \\
		\cadreblanc{10}
		\item $a_0 = \dfrac{1}{T}\displaystyle\int_0^{T} f(t)\:\text{d}t$. Justifier que $a_0 = 9,6$.\\
		\cadreblanc{20}
	\end{enumerate}
\item La valeur efficace de $U$, notée $U_{\text{eff}}$, vérifie : $U^2_{\text{eff}} = \dfrac{1}{T}\displaystyle\int_0^{T} [f(t)]^2\: \text{d}t$. 

Montrer que : $U_{\text{eff}} = 10,7$~V.\\
\cadreblanc{20}
\item Le signal correspondant à la tension $U$ est envoyé sur une ligne moyenne tension transportant une tension efficace de $M = 20 000$~V.

Le taux de distorsion harmonique par rapport au fondamental, noté $T_F$, est donné par la formule suivante:

\[T_F = \sqrt{\dfrac{U^2_{\text{eff}} - a_0^2}{M^2}}.\]

On considère qu'un CPL n'a pas d'incidence sur le réseau si $T_F$ est inférieur à 0,1\,\%. 

Le CPL étudié dans la partie A a-t-il une incidence sur le réseau ?\\
\cadreblanc{30}
\end{enumerate}


\newpage 


\fcolorbox[gray]{0}{0.9}{\textbf{Extrait - Brevet de technicien supérieur septembre 2020 - groupement A}}\\[2mm]
Le moteur triphasé d'une voiture électrique est alimenté par un onduleur pour modifier sa vitesse. 

La tension simple à la sortie de cet onduleur dépend de $\omega t$  où $\omega$ est la pulsation (en rad/s) et $t$ le temps (en s). On l'exprime sous la forme $v(\omega t)$, où $v$ est une fonction paire, périodique de période $T = 2\pi$, et définie pour tout réel $x$ appartenant à l'intervalle $[0~;~\pi]$ par :

\[\left\{\begin{array}{l c l c l}
v(x) &=& 300 &\text{si} &0 \leqslant x < \frac{\pi}{6} \\[7pt]
v(x) &=& 150 &\text{si} & \frac{\pi}{6} \leqslant  x \leqslant \frac{\pi}{2}\\[7pt]
v(x) &=& -150 &\text{si} & \frac{\pi}{2} < x < \frac{5\pi}{6}\\[7pt]
v(x) &=& -300 &\text{si} &\frac{5\pi}{6} \leqslant x \leqslant \pi
\end{array}\right.\]

\medskip

\begin{enumerate}
\item Compléter la représentation graphique de la fonction $v$ sur l'intervalle $\left[- \dfrac{3\pi}{2}~;~\dfrac{3\pi}{2}\right]$.\\
\begin{center}
% 1 unité x = π/6 → xmin=-9 (=-3π/2), xmax=9 (=3π/2)
% 1 unité y = 50V → ymin=-7 (=-350V), ymax=7 (=350V)
\begin{tikzpicture}[x=0.7cm, y=0.38cm]
\def\xmin{-9.5} \def\xmax{9.5} \def\ymin{-7.5} \def\ymax{7.5}
% Grilles (xstep en unités = π/6, ystep=50V=1unité)
\draw[xstep=1,ystep=1,very thin,orange!25] (\xmin,\ymin) grid (\xmax,\ymax);
\draw[xstep=3,ystep=2,thin,orange!55] (\xmin,\ymin) grid (\xmax,\ymax);
% Axes
\draw[->,>=stealth',thick,blue!70!black] (\xmin,0)--(\xmax,0) node[right]{\small$x$};
\draw[->,>=stealth',thick,blue!70!black] (0,\ymin)--(0,\ymax) node[above]{\small$v$};
% Graduations x : -9=−3π/2, -6=−π, -3=−π/2, 3=π/2, 6=π, 9=3π/2
\foreach \n/\lab in {-9/{-\frac{3\pi}{2}},-6/{-\pi},-3/{-\frac{\pi}{2}},3/{\frac{\pi}{2}},6/{\pi},9/{\frac{3\pi}{2}}}
	\draw[blue!70!black] (\n,3pt)--(\n,-5pt) node[below]{\small$\lab$};
% Graduations y : valeurs 150 et 300 (et opposées)
\foreach \y/\lab in {-6/{-300},-3/{-150},3/{150},6/{300}}
	\draw[blue!70!black] (2pt,\y)--(-5pt,\y) node[left]{\small$\lab$};
% Signal sur [0, π] (6 unités x, en unités de π/6) :
%   [0, π/6[ : v=300  → x∈[0,1[
%   [π/6, π/2] : v=150 → x∈[1,3]
%   ]π/2, 5π/6[ : v=-150 → x∈[3,5[
%   [5π/6, π] : v=-300  → x∈[5,6]
% Signal tracé partiellement — le début [0,π] est donné
\draw[very thick,red!70!black]
	(0,6)--(1,6)--(1,3)--(3,3)--(3,-3)--(5,-3)--(5,-6)--(6,-6);
% Traits verticaux
\draw[very thick,red!70!black] (0,-7)--(0,7);
% NOTE : la partie [-3π/2 ; 0] est à compléter par l'élève (parité + périodicité)
% On trace seulement un léger repère pointillé pour guider
\draw[dashed,blue!30] (-9,0)--(0,0);
\draw[black!60,thick] (\xmin,\ymin) rectangle (\xmax,\ymax);
\end{tikzpicture}
\end{center}

On admet que la fonction $v$ est développable en série de Fourier et que, pour tout réel $x$, on a:

\[v(x) = a_0 + \displaystyle\sum_{n=1}^{+ \infty} \left(a_n \cos (nx) + b_n \sin (nx) \right).\]

\item Justifier que $b_n = 0$ pour tout entier $n\geqslant 1$.\\
\cadreblanc{10}
\item On rappelle que : $a_0 = \dfrac{2}{T}\displaystyle\int_0^{\frac{T}{2}} v(x)\:\text{d}x$.


 Montrer que : $a_0 = 0$.\\
 \cadreblanc{20}
 
\item On donne, pour tout entier $n \geqslant 1$ :\: $a_n = \dfrac{4}{T}\displaystyle\int_0^{\frac{T}{2}}v(x)\cos(nx)\:\text{d}x$.
	\begin{enumerate}
		\item On pose : $I_n =  \displaystyle\int_0^{\frac{\pi}{6}} 300\, \cos(nx)\:\text{d}x$ et $J_n =\displaystyle\int_{\frac{5\pi}{6}}^{\pi} \left[- 300 \cos(nx)\right]\:\text{d}x$
		
Calculer $I_n$ et $J_n$ en fonction de $n$.\\
\cadreblanc{40}

		\item Écrire les intégrales qu'il reste à calculer pour obtenir l'expression de $a_n$ en fonction de $n$. 
		
\textbf{On ne demande pas de les calculer.}\\
\cadreblanc{40}

	\end{enumerate}
\item On admet que, pour tout entier $n \geqslant 1$, on a : $a_n = \dfrac{300}{n\pi} \left[\sin\left(\dfrac{n\pi}{6}\right) + 2 \sin \left(\dfrac{n\pi}{2}\right) + \sin \left(\dfrac{5n \pi}{6}\right)\right]$.

Compléter le tableau ci-dessous :
\begin{center}
\begin{tabularx}{\linewidth}{|m{4cm}|*{8}{>{\centering \arraybackslash}X|}}\hline
Valeurs de $n$			& 0&1		&2	&3	&4	&5	&6	&7\\ \hline
Valeurs approchées de $a_n$
 arrondies au dixième.	&0 &286,5 	&0 	&0	&	&	&	&\\ \hline
\end{tabularx}
\end{center}


\item  On évalue la pollution harmonique de deux façons : de manière qualitative, en observant la présence d'harmoniques sur le spectre des amplitudes, et de manière quantitative en calculant le taux de distorsion harmonique par rapport au fondamental.
	\begin{enumerate}
		\item Compléter sur le graphique ci-dessous le spectre des amplitudes $A_n$ associé à la fonction $v$ pour $n$ allant de $0$ à $7$.\\
		On rappelle que pour tout entier $n$ supérieur ou égal à $1$ : $A_n = \sqrt{a_n^2  + b_n^2}$.
		\begin{center}
		\begin{tikzpicture}[x=1.6cm, y=0.013cm]
\def\xmin{-0.5} \def\xmax{7.8} \def\ymin{-20} \def\ymax{320}
% Grilles
\draw[xstep=1,ystep=50,very thin,orange!25] (\xmin,0) grid (\xmax,\ymax);
\draw[xstep=1,ystep=100,thin,orange!55] (\xmin,0) grid (\xmax,\ymax);
% Axes
\draw[->,>=stealth',thick,blue!70!black] (\xmin,0)--(\xmax,0) node[right]{\small$n$};
\draw[->,>=stealth',thick,blue!70!black] (0,\ymin)--(0,\ymax) node[above]{\small$A_n$~(V)};
% Graduations x
\foreach \n in {0,1,2,3,4,5,6,7}
	\draw[blue!70!black] (\n,3pt)--(\n,-8pt) node[below]{\small$\n$};
% Graduations y
\foreach \y in {50,100,150,200,250,300}
	\draw[blue!70!black] (2pt,\y)--(-6pt,\y) node[left]{\small$\y$};
% Seule barre pré-tracée : A_1 = 900/π ≈ 286.5 V
\fill[red!60!white] (0.7,0) rectangle (1.3,286.5);
\draw[red!70!black,thick] (0.7,0) rectangle (1.3,286.5);
\node[above,red!70!black,font=\small] at (1,286.5) {$\approx 287$};
% Cadres vides pour n=0,2,3,4,5,6,7
\foreach \n in {0,2,3,4,5,6,7}
	\draw[blue!50!black,thick,dashed] (\n-0.3,0) rectangle (\n+0.3,20);
\draw[black!60,thick] (\xmin,\ymin) rectangle (\xmax,\ymax);
\end{tikzpicture}
		\end{center}
		

		\item On s'intéresse au taux de distorsion harmonique par rapport au fondamental.
		
On obtient une valeur approchée significative de ce taux en calculant le nombre T donné par :

\[\text{T} = \dfrac{\sqrt{a_2^2 + a_3^2 + a_4^2 + a_5^2 + a_6^2 + a_7^2}}{a_1}.\]

Sachant que ce taux doit être inférieur à 10\,\%, la commande de l'onduleur du moteur est-elle adaptée ?\\
\cadreblanc{25}
	\end{enumerate}
\end{enumerate}


\setlength{\parindent}{0cm}


\bigskip

%%%%%%%%%%%%%%%%%%%%%%%%%%%%%%%%%%%%%%%%%
\newpage
%%%%DEVOIR SURVEILLE%%%%%%%%%%%%%%%%%%%%%

%!TEX encoding = UTF-8 Unicode



\end{document}
